\QHeading{第1题　在既有平行线段又有垂直线段的图形下面的（\quad）里画“$\checkmark$”。}

\par\noindent\makebox[\linewidth][c]{
\begin{tikzpicture}[x=1bp,y=-1bp,color=black,draw=black,text=black]
\begin{scope}\clip (56.0,131.0) rectangle (469.0,221.33);
\hdhead{94}{总复习}
\hdsection{1}{150}{在既有平行线段又有垂直线段的图形下面的（\quad）里画“$\checkmark$”。}
\draw[line width=.7pt] (148.000,168.000)--(170.000,136.667)--(220.000,136.667)--(197.333,168.000)--cycle;
\draw[line width=.7pt] (242.667,136.667)--(308.667,136.667)--(308.667,168.000)--(242.667,168.000)--cycle;
\draw[line width=.7pt] (334.667,136.667)--(372.000,136.667)--(385.333,168.000)--(334.667,168.000)--cycle;
\draw[line width=.7pt] (406.667,136.667)--(406.667,168.000)--(462.667,168.000)--cycle;
\q{236}{279}{110}{12}{（\hspace{30pt}）}
\q{374}{279}{110}{12}{（\hspace{30pt}）}
\q{515}{279}{110}{12}{（\hspace{30pt}）}
\q{615}{279}{110}{12}{（\hspace{30pt}）}
\hdsection{2}{335}{数一数，填一填。}
\q{102}{380}{42}{13}{①}
\q{403}{380}{42}{13}{②}
\draw[line width=.7pt] (89.333,302.000)--(116.667,256.000)--(143.333,302.000)--cycle;
\draw[line width=.7pt] (103.000,279.000)--(130.000,279.000)--(116.667,302.000)--cycle;
\draw[line width=.7pt] (290.000,302.000)--(316.667,256.000)--(374.667,256.000)--(402.000,302.000)--cycle;
\draw[line width=.7pt] (316.667,256.000)--(343.333,302.000);
\draw[line width=.7pt] (303.333,279.000)--(330.000,279.000)--(330.000,302.000);
\q{215}{387}{196}{12}{（\hspace{25pt}）个三角形}
\q{215}{431}{196}{12}{（\hspace{25pt}）个梯形}
\q{603}{387}{230}{12}{（\hspace{25pt}）个梯形}
\q{603}{431}{250}{12}{（\hspace{25pt}）个平行四边形}
\hdsection{3}{490}{分一分。（填序号。）}
\draw[line width=.7pt] (139.333,396.000)--(152.667,363.333)--(191.333,363.333)--(177.333,396.000)--cycle;
\draw[line width=.7pt] (245.333,395.333)--(260.667,364.667)--(275.333,395.333)--cycle;
\draw[line width=.7pt] (326.000,363.333)--(326.000,396.000)--(384.667,396.000)--cycle;
\draw[line width=.7pt] (416.667,394.000)--(432.667,364.667)--(466.667,364.667)--(482.667,394.000)--cycle;
\draw[line width=.7pt] (146.000,468.000)--(164.667,434.000)--(184.667,468.000)--cycle;
\draw[line width=.7pt] (238.667,434.667)--(282.667,434.667)--(282.667,468.000)--(238.667,468.000)--cycle;
\draw[line width=.7pt] (338.000,448.000)--(354.667,434.667)--(372.000,448.000)--(366.667,468.000)--(344.000,468.000)--cycle;
\draw[line width=.7pt] (439.333,467.333)--(444.000,434.000)--(460.667,434.000)--(456.000,467.333)--cycle;
\node[inner sep=0pt,align=center,font=\normalsize,] at (165.333,414.000) {①};
\node[inner sep=0pt,align=center,font=\normalsize,] at (260.667,408.667) {②};
\node[inner sep=0pt,align=center,font=\normalsize,] at (355.333,413.333) {③};
\node[inner sep=0pt,align=center,font=\normalsize,] at (450.000,408.667) {④};
\node[inner sep=0pt,align=center,font=\normalsize,] at (165.333,485.333) {⑤};
\node[inner sep=0pt,align=center,font=\normalsize,] at (260.667,485.333) {⑥};
\node[inner sep=0pt,align=center,font=\normalsize,] at (355.333,485.333) {⑦};
\node[inner sep=0pt,align=center,font=\normalsize,] at (450.000,481.333) {⑧};
\q{104}{769}{241}{12}{三角形：（\hspace{46pt}）}
\q{329}{769}{292}{12}{平行四边形：（\hspace{50pt}）}
\q{626}{769}{200}{12}{梯形：（\hspace{47pt}）}
\hdsection{4}{827}{在下面的点子图中画一画。}
\q{103}{870}{722}{12}{① 画出有一个角是直角的等腰三角形。}
\q{103}{914}{722}{12}{② 画一个梯形，并画出它的一条高，标出上底和下底。}
\q{103}{957}{722}{12}{③ 画一个平行四边形，并画出它的一条高。}
\fill (75.333,673.333) circle (.7bp);
\fill (87.333,673.333) circle (.7bp);
\fill (99.333,673.333) circle (.7bp);
\fill (111.333,673.333) circle (.7bp);
\fill (123.333,673.333) circle (.7bp);
\fill (135.333,673.333) circle (.7bp);
\fill (147.333,673.333) circle (.7bp);
\fill (159.333,673.333) circle (.7bp);
\fill (171.333,673.333) circle (.7bp);
\fill (183.333,673.333) circle (.7bp);
\fill (195.333,673.333) circle (.7bp);
\fill (207.333,673.333) circle (.7bp);
\fill (219.333,673.333) circle (.7bp);
\fill (231.333,673.333) circle (.7bp);
\fill (243.333,673.333) circle (.7bp);
\fill (255.333,673.333) circle (.7bp);
\fill (267.333,673.333) circle (.7bp);
\fill (279.333,673.333) circle (.7bp);
\fill (291.333,673.333) circle (.7bp);
\fill (303.333,673.333) circle (.7bp);
\fill (315.333,673.333) circle (.7bp);
\fill (327.333,673.333) circle (.7bp);
\fill (339.333,673.333) circle (.7bp);
\fill (351.333,673.333) circle (.7bp);
\fill (363.333,673.333) circle (.7bp);
\fill (375.333,673.333) circle (.7bp);
\fill (387.333,673.333) circle (.7bp);
\fill (399.333,673.333) circle (.7bp);
\fill (411.333,673.333) circle (.7bp);
\fill (423.333,673.333) circle (.7bp);
\fill (435.333,673.333) circle (.7bp);
\fill (447.333,673.333) circle (.7bp);
\fill (459.333,673.333) circle (.7bp);
\fill (471.333,673.333) circle (.7bp);
\fill (483.333,673.333) circle (.7bp);
\fill (495.333,673.333) circle (.7bp);
\fill (507.333,673.333) circle (.7bp);
\fill (519.333,673.333) circle (.7bp);
\fill (531.333,673.333) circle (.7bp);
\fill (543.333,673.333) circle (.7bp);
\fill (75.333,685.333) circle (.7bp);
\fill (87.333,685.333) circle (.7bp);
\fill (99.333,685.333) circle (.7bp);
\fill (111.333,685.333) circle (.7bp);
\fill (123.333,685.333) circle (.7bp);
\fill (135.333,685.333) circle (.7bp);
\fill (147.333,685.333) circle (.7bp);
\fill (159.333,685.333) circle (.7bp);
\fill (171.333,685.333) circle (.7bp);
\fill (183.333,685.333) circle (.7bp);
\fill (195.333,685.333) circle (.7bp);
\fill (207.333,685.333) circle (.7bp);
\fill (219.333,685.333) circle (.7bp);
\fill (231.333,685.333) circle (.7bp);
\fill (243.333,685.333) circle (.7bp);
\fill (255.333,685.333) circle (.7bp);
\fill (267.333,685.333) circle (.7bp);
\fill (279.333,685.333) circle (.7bp);
\fill (291.333,685.333) circle (.7bp);
\fill (303.333,685.333) circle (.7bp);
\fill (315.333,685.333) circle (.7bp);
\fill (327.333,685.333) circle (.7bp);
\fill (339.333,685.333) circle (.7bp);
\fill (351.333,685.333) circle (.7bp);
\fill (363.333,685.333) circle (.7bp);
\fill (375.333,685.333) circle (.7bp);
\fill (387.333,685.333) circle (.7bp);
\fill (399.333,685.333) circle (.7bp);
\fill (411.333,685.333) circle (.7bp);
\fill (423.333,685.333) circle (.7bp);
\fill (435.333,685.333) circle (.7bp);
\fill (447.333,685.333) circle (.7bp);
\fill (459.333,685.333) circle (.7bp);
\fill (471.333,685.333) circle (.7bp);
\fill (483.333,685.333) circle (.7bp);
\fill (495.333,685.333) circle (.7bp);
\fill (507.333,685.333) circle (.7bp);
\fill (519.333,685.333) circle (.7bp);
\fill (531.333,685.333) circle (.7bp);
\fill (543.333,685.333) circle (.7bp);
\fill (75.333,697.333) circle (.7bp);
\fill (87.333,697.333) circle (.7bp);
\fill (99.333,697.333) circle (.7bp);
\fill (111.333,697.333) circle (.7bp);
\fill (123.333,697.333) circle (.7bp);
\fill (135.333,697.333) circle (.7bp);
\fill (147.333,697.333) circle (.7bp);
\fill (159.333,697.333) circle (.7bp);
\fill (171.333,697.333) circle (.7bp);
\fill (183.333,697.333) circle (.7bp);
\fill (195.333,697.333) circle (.7bp);
\fill (207.333,697.333) circle (.7bp);
\fill (219.333,697.333) circle (.7bp);
\fill (231.333,697.333) circle (.7bp);
\fill (243.333,697.333) circle (.7bp);
\fill (255.333,697.333) circle (.7bp);
\fill (267.333,697.333) circle (.7bp);
\fill (279.333,697.333) circle (.7bp);
\fill (291.333,697.333) circle (.7bp);
\fill (303.333,697.333) circle (.7bp);
\fill (315.333,697.333) circle (.7bp);
\fill (327.333,697.333) circle (.7bp);
\fill (339.333,697.333) circle (.7bp);
\fill (351.333,697.333) circle (.7bp);
\fill (363.333,697.333) circle (.7bp);
\fill (375.333,697.333) circle (.7bp);
\fill (387.333,697.333) circle (.7bp);
\fill (399.333,697.333) circle (.7bp);
\fill (411.333,697.333) circle (.7bp);
\fill (423.333,697.333) circle (.7bp);
\fill (435.333,697.333) circle (.7bp);
\fill (447.333,697.333) circle (.7bp);
\fill (459.333,697.333) circle (.7bp);
\fill (471.333,697.333) circle (.7bp);
\fill (483.333,697.333) circle (.7bp);
\fill (495.333,697.333) circle (.7bp);
\fill (507.333,697.333) circle (.7bp);
\fill (519.333,697.333) circle (.7bp);
\fill (531.333,697.333) circle (.7bp);
\fill (543.333,697.333) circle (.7bp);
\fill (75.333,709.333) circle (.7bp);
\fill (87.333,709.333) circle (.7bp);
\fill (99.333,709.333) circle (.7bp);
\fill (111.333,709.333) circle (.7bp);
\fill (123.333,709.333) circle (.7bp);
\fill (135.333,709.333) circle (.7bp);
\fill (147.333,709.333) circle (.7bp);
\fill (159.333,709.333) circle (.7bp);
\fill (171.333,709.333) circle (.7bp);
\fill (183.333,709.333) circle (.7bp);
\fill (195.333,709.333) circle (.7bp);
\fill (207.333,709.333) circle (.7bp);
\fill (219.333,709.333) circle (.7bp);
\fill (231.333,709.333) circle (.7bp);
\fill (243.333,709.333) circle (.7bp);
\fill (255.333,709.333) circle (.7bp);
\fill (267.333,709.333) circle (.7bp);
\fill (279.333,709.333) circle (.7bp);
\fill (291.333,709.333) circle (.7bp);
\fill (303.333,709.333) circle (.7bp);
\fill (315.333,709.333) circle (.7bp);
\fill (327.333,709.333) circle (.7bp);
\fill (339.333,709.333) circle (.7bp);
\fill (351.333,709.333) circle (.7bp);
\fill (363.333,709.333) circle (.7bp);
\fill (375.333,709.333) circle (.7bp);
\fill (387.333,709.333) circle (.7bp);
\fill (399.333,709.333) circle (.7bp);
\fill (411.333,709.333) circle (.7bp);
\fill (423.333,709.333) circle (.7bp);
\fill (435.333,709.333) circle (.7bp);
\fill (447.333,709.333) circle (.7bp);
\fill (459.333,709.333) circle (.7bp);
\fill (471.333,709.333) circle (.7bp);
\fill (483.333,709.333) circle (.7bp);
\fill (495.333,709.333) circle (.7bp);
\fill (507.333,709.333) circle (.7bp);
\fill (519.333,709.333) circle (.7bp);
\fill (531.333,709.333) circle (.7bp);
\fill (543.333,709.333) circle (.7bp);
\fill (75.333,721.333) circle (.7bp);
\fill (87.333,721.333) circle (.7bp);
\fill (99.333,721.333) circle (.7bp);
\fill (111.333,721.333) circle (.7bp);
\fill (123.333,721.333) circle (.7bp);
\fill (135.333,721.333) circle (.7bp);
\fill (147.333,721.333) circle (.7bp);
\fill (159.333,721.333) circle (.7bp);
\fill (171.333,721.333) circle (.7bp);
\fill (183.333,721.333) circle (.7bp);
\fill (195.333,721.333) circle (.7bp);
\fill (207.333,721.333) circle (.7bp);
\fill (219.333,721.333) circle (.7bp);
\fill (231.333,721.333) circle (.7bp);
\fill (243.333,721.333) circle (.7bp);
\fill (255.333,721.333) circle (.7bp);
\fill (267.333,721.333) circle (.7bp);
\fill (279.333,721.333) circle (.7bp);
\fill (291.333,721.333) circle (.7bp);
\fill (303.333,721.333) circle (.7bp);
\fill (315.333,721.333) circle (.7bp);
\fill (327.333,721.333) circle (.7bp);
\fill (339.333,721.333) circle (.7bp);
\fill (351.333,721.333) circle (.7bp);
\fill (363.333,721.333) circle (.7bp);
\fill (375.333,721.333) circle (.7bp);
\fill (387.333,721.333) circle (.7bp);
\fill (399.333,721.333) circle (.7bp);
\fill (411.333,721.333) circle (.7bp);
\fill (423.333,721.333) circle (.7bp);
\fill (435.333,721.333) circle (.7bp);
\fill (447.333,721.333) circle (.7bp);
\fill (459.333,721.333) circle (.7bp);
\fill (471.333,721.333) circle (.7bp);
\fill (483.333,721.333) circle (.7bp);
\fill (495.333,721.333) circle (.7bp);
\fill (507.333,721.333) circle (.7bp);
\fill (519.333,721.333) circle (.7bp);
\fill (531.333,721.333) circle (.7bp);
\fill (543.333,721.333) circle (.7bp);
\fill (75.333,733.333) circle (.7bp);
\fill (87.333,733.333) circle (.7bp);
\fill (99.333,733.333) circle (.7bp);
\fill (111.333,733.333) circle (.7bp);
\fill (123.333,733.333) circle (.7bp);
\fill (135.333,733.333) circle (.7bp);
\fill (147.333,733.333) circle (.7bp);
\fill (159.333,733.333) circle (.7bp);
\fill (171.333,733.333) circle (.7bp);
\fill (183.333,733.333) circle (.7bp);
\fill (195.333,733.333) circle (.7bp);
\fill (207.333,733.333) circle (.7bp);
\fill (219.333,733.333) circle (.7bp);
\fill (231.333,733.333) circle (.7bp);
\fill (243.333,733.333) circle (.7bp);
\fill (255.333,733.333) circle (.7bp);
\fill (267.333,733.333) circle (.7bp);
\fill (279.333,733.333) circle (.7bp);
\fill (291.333,733.333) circle (.7bp);
\fill (303.333,733.333) circle (.7bp);
\fill (315.333,733.333) circle (.7bp);
\fill (327.333,733.333) circle (.7bp);
\fill (339.333,733.333) circle (.7bp);
\fill (351.333,733.333) circle (.7bp);
\fill (363.333,733.333) circle (.7bp);
\fill (375.333,733.333) circle (.7bp);
\fill (387.333,733.333) circle (.7bp);
\fill (399.333,733.333) circle (.7bp);
\fill (411.333,733.333) circle (.7bp);
\fill (423.333,733.333) circle (.7bp);
\fill (435.333,733.333) circle (.7bp);
\fill (447.333,733.333) circle (.7bp);
\fill (459.333,733.333) circle (.7bp);
\fill (471.333,733.333) circle (.7bp);
\fill (483.333,733.333) circle (.7bp);
\fill (495.333,733.333) circle (.7bp);
\fill (507.333,733.333) circle (.7bp);
\fill (519.333,733.333) circle (.7bp);
\fill (531.333,733.333) circle (.7bp);
\fill (543.333,733.333) circle (.7bp);
\fill (75.333,745.333) circle (.7bp);
\fill (87.333,745.333) circle (.7bp);
\fill (99.333,745.333) circle (.7bp);
\fill (111.333,745.333) circle (.7bp);
\fill (123.333,745.333) circle (.7bp);
\fill (135.333,745.333) circle (.7bp);
\fill (147.333,745.333) circle (.7bp);
\fill (159.333,745.333) circle (.7bp);
\fill (171.333,745.333) circle (.7bp);
\fill (183.333,745.333) circle (.7bp);
\fill (195.333,745.333) circle (.7bp);
\fill (207.333,745.333) circle (.7bp);
\fill (219.333,745.333) circle (.7bp);
\fill (231.333,745.333) circle (.7bp);
\fill (243.333,745.333) circle (.7bp);
\fill (255.333,745.333) circle (.7bp);
\fill (267.333,745.333) circle (.7bp);
\fill (279.333,745.333) circle (.7bp);
\fill (291.333,745.333) circle (.7bp);
\fill (303.333,745.333) circle (.7bp);
\fill (315.333,745.333) circle (.7bp);
\fill (327.333,745.333) circle (.7bp);
\fill (339.333,745.333) circle (.7bp);
\fill (351.333,745.333) circle (.7bp);
\fill (363.333,745.333) circle (.7bp);
\fill (375.333,745.333) circle (.7bp);
\fill (387.333,745.333) circle (.7bp);
\fill (399.333,745.333) circle (.7bp);
\fill (411.333,745.333) circle (.7bp);
\fill (423.333,745.333) circle (.7bp);
\fill (435.333,745.333) circle (.7bp);
\fill (447.333,745.333) circle (.7bp);
\fill (459.333,745.333) circle (.7bp);
\fill (471.333,745.333) circle (.7bp);
\fill (483.333,745.333) circle (.7bp);
\fill (495.333,745.333) circle (.7bp);
\fill (507.333,745.333) circle (.7bp);
\fill (519.333,745.333) circle (.7bp);
\fill (531.333,745.333) circle (.7bp);
\fill (543.333,745.333) circle (.7bp);
\fill (75.333,757.333) circle (.7bp);
\fill (87.333,757.333) circle (.7bp);
\fill (99.333,757.333) circle (.7bp);
\fill (111.333,757.333) circle (.7bp);
\fill (123.333,757.333) circle (.7bp);
\fill (135.333,757.333) circle (.7bp);
\fill (147.333,757.333) circle (.7bp);
\fill (159.333,757.333) circle (.7bp);
\fill (171.333,757.333) circle (.7bp);
\fill (183.333,757.333) circle (.7bp);
\fill (195.333,757.333) circle (.7bp);
\fill (207.333,757.333) circle (.7bp);
\fill (219.333,757.333) circle (.7bp);
\fill (231.333,757.333) circle (.7bp);
\fill (243.333,757.333) circle (.7bp);
\fill (255.333,757.333) circle (.7bp);
\fill (267.333,757.333) circle (.7bp);
\fill (279.333,757.333) circle (.7bp);
\fill (291.333,757.333) circle (.7bp);
\fill (303.333,757.333) circle (.7bp);
\fill (315.333,757.333) circle (.7bp);
\fill (327.333,757.333) circle (.7bp);
\fill (339.333,757.333) circle (.7bp);
\fill (351.333,757.333) circle (.7bp);
\fill (363.333,757.333) circle (.7bp);
\fill (375.333,757.333) circle (.7bp);
\fill (387.333,757.333) circle (.7bp);
\fill (399.333,757.333) circle (.7bp);
\fill (411.333,757.333) circle (.7bp);
\fill (423.333,757.333) circle (.7bp);
\fill (435.333,757.333) circle (.7bp);
\fill (447.333,757.333) circle (.7bp);
\fill (459.333,757.333) circle (.7bp);
\fill (471.333,757.333) circle (.7bp);
\fill (483.333,757.333) circle (.7bp);
\fill (495.333,757.333) circle (.7bp);
\fill (507.333,757.333) circle (.7bp);
\fill (519.333,757.333) circle (.7bp);
\fill (531.333,757.333) circle (.7bp);
\fill (543.333,757.333) circle (.7bp);
\fill (75.333,769.333) circle (.7bp);
\fill (87.333,769.333) circle (.7bp);
\fill (99.333,769.333) circle (.7bp);
\fill (111.333,769.333) circle (.7bp);
\fill (123.333,769.333) circle (.7bp);
\fill (135.333,769.333) circle (.7bp);
\fill (147.333,769.333) circle (.7bp);
\fill (159.333,769.333) circle (.7bp);
\fill (171.333,769.333) circle (.7bp);
\fill (183.333,769.333) circle (.7bp);
\fill (195.333,769.333) circle (.7bp);
\fill (207.333,769.333) circle (.7bp);
\fill (219.333,769.333) circle (.7bp);
\fill (231.333,769.333) circle (.7bp);
\fill (243.333,769.333) circle (.7bp);
\fill (255.333,769.333) circle (.7bp);
\fill (267.333,769.333) circle (.7bp);
\fill (279.333,769.333) circle (.7bp);
\fill (291.333,769.333) circle (.7bp);
\fill (303.333,769.333) circle (.7bp);
\fill (315.333,769.333) circle (.7bp);
\fill (327.333,769.333) circle (.7bp);
\fill (339.333,769.333) circle (.7bp);
\fill (351.333,769.333) circle (.7bp);
\fill (363.333,769.333) circle (.7bp);
\fill (375.333,769.333) circle (.7bp);
\fill (387.333,769.333) circle (.7bp);
\fill (399.333,769.333) circle (.7bp);
\fill (411.333,769.333) circle (.7bp);
\fill (423.333,769.333) circle (.7bp);
\fill (435.333,769.333) circle (.7bp);
\fill (447.333,769.333) circle (.7bp);
\fill (459.333,769.333) circle (.7bp);
\fill (471.333,769.333) circle (.7bp);
\fill (483.333,769.333) circle (.7bp);
\fill (495.333,769.333) circle (.7bp);
\fill (507.333,769.333) circle (.7bp);
\fill (519.333,769.333) circle (.7bp);
\fill (531.333,769.333) circle (.7bp);
\fill (543.333,769.333) circle (.7bp);

\begin{scope}[color=answerred,draw=answerred,text=answerred]% 高清文字版专用填空位置；原页答案另存于 ../overlays。
\ans{268.667}{186.667}{40.000}{14}{$\checkmark$}
\ans{352.000}{186.667}{40.000}{14}{$\checkmark$}
\ans{162.000}{263.333}{32.000}{12}{$5$}
\ans{162.000}{292.667}{32.000}{12}{$3$}
\ans{420.667}{263.333}{32.000}{12}{$3$}
\ans{420.667}{292.667}{32.000}{12}{$1$}
\ans{129.613}{512.343}{66.000}{11}{②③⑤}
\ans{305.517}{512.343}{70.000}{11}{①⑥⑧}
\ans{477.115}{512.343}{43.333}{11}{④}
\begin{scope}[scale=0.6666667]\draw[answerred,line width=1pt] (149,1136)--(149,1064)--(221,1136)--cycle;\draw (149,1125)--(160,1125)--(160,1136);\draw (144,1100)--(154,1100) (185,1131)--(185,1141);\node[answerred,above] at (185,1050) {①};\draw[answerred,line width=1pt] (365,1136)--(401,1064)--(473,1064)--(509,1136)--cycle;\draw[answerred,dashed] (401,1064)--(401,1136);\draw (401,1125)--(412,1125)--(412,1136);\node[above] at (437,1061) {上底};\node[below] at (437,1136) {下底};\node[left] at (401,1100) {高};\node[above] at (365,1050) {②};\draw[answerred,line width=1pt] (617,1136)--(653,1064)--(761,1064)--(725,1136)--cycle;\draw[answerred,dashed] (653,1064)--(653,1136);\draw (653,1125)--(664,1125)--(664,1136);\node[left] at (653,1100) {高};\node[above] at (616,1050) {③};\end{scope}
\end{scope}
\end{scope}\pgfresetboundingbox\path[use as bounding box] (56.0,131.0) rectangle (469.0,221.33);
\end{tikzpicture}}\par
\par\noindent




\par



\par\noindent

\QHeading{第2题　数一数，填一填。}

\par\noindent\makebox[\linewidth][c]{
\begin{tikzpicture}[x=1bp,y=-1bp,color=black,draw=black,text=black]
\begin{scope}\clip (56.0,248.0) rectangle (527.0,324.67);
\hdhead{94}{总复习}
\hdsection{1}{150}{在既有平行线段又有垂直线段的图形下面的（\quad）里画“$\checkmark$”。}
\draw[line width=.7pt] (148.000,168.000)--(170.000,136.667)--(220.000,136.667)--(197.333,168.000)--cycle;
\draw[line width=.7pt] (242.667,136.667)--(308.667,136.667)--(308.667,168.000)--(242.667,168.000)--cycle;
\draw[line width=.7pt] (334.667,136.667)--(372.000,136.667)--(385.333,168.000)--(334.667,168.000)--cycle;
\draw[line width=.7pt] (406.667,136.667)--(406.667,168.000)--(462.667,168.000)--cycle;
\q{236}{279}{110}{12}{（\hspace{30pt}）}
\q{374}{279}{110}{12}{（\hspace{30pt}）}
\q{515}{279}{110}{12}{（\hspace{30pt}）}
\q{615}{279}{110}{12}{（\hspace{30pt}）}
\hdsection{2}{335}{数一数，填一填。}
\q{102}{380}{42}{13}{①}
\q{403}{380}{42}{13}{②}
\draw[line width=.7pt] (89.333,302.000)--(116.667,256.000)--(143.333,302.000)--cycle;
\draw[line width=.7pt] (103.000,279.000)--(130.000,279.000)--(116.667,302.000)--cycle;
\draw[line width=.7pt] (290.000,302.000)--(316.667,256.000)--(374.667,256.000)--(402.000,302.000)--cycle;
\draw[line width=.7pt] (316.667,256.000)--(343.333,302.000);
\draw[line width=.7pt] (303.333,279.000)--(330.000,279.000)--(330.000,302.000);
\q{215}{387}{196}{12}{（\hspace{25pt}）个三角形}
\q{215}{431}{196}{12}{（\hspace{25pt}）个梯形}
\q{603}{387}{230}{12}{（\hspace{25pt}）个梯形}
\q{603}{431}{250}{12}{（\hspace{25pt}）个平行四边形}
\hdsection{3}{490}{分一分。（填序号。）}
\draw[line width=.7pt] (139.333,396.000)--(152.667,363.333)--(191.333,363.333)--(177.333,396.000)--cycle;
\draw[line width=.7pt] (245.333,395.333)--(260.667,364.667)--(275.333,395.333)--cycle;
\draw[line width=.7pt] (326.000,363.333)--(326.000,396.000)--(384.667,396.000)--cycle;
\draw[line width=.7pt] (416.667,394.000)--(432.667,364.667)--(466.667,364.667)--(482.667,394.000)--cycle;
\draw[line width=.7pt] (146.000,468.000)--(164.667,434.000)--(184.667,468.000)--cycle;
\draw[line width=.7pt] (238.667,434.667)--(282.667,434.667)--(282.667,468.000)--(238.667,468.000)--cycle;
\draw[line width=.7pt] (338.000,448.000)--(354.667,434.667)--(372.000,448.000)--(366.667,468.000)--(344.000,468.000)--cycle;
\draw[line width=.7pt] (439.333,467.333)--(444.000,434.000)--(460.667,434.000)--(456.000,467.333)--cycle;
\node[inner sep=0pt,align=center,font=\normalsize,] at (165.333,414.000) {①};
\node[inner sep=0pt,align=center,font=\normalsize,] at (260.667,408.667) {②};
\node[inner sep=0pt,align=center,font=\normalsize,] at (355.333,413.333) {③};
\node[inner sep=0pt,align=center,font=\normalsize,] at (450.000,408.667) {④};
\node[inner sep=0pt,align=center,font=\normalsize,] at (165.333,485.333) {⑤};
\node[inner sep=0pt,align=center,font=\normalsize,] at (260.667,485.333) {⑥};
\node[inner sep=0pt,align=center,font=\normalsize,] at (355.333,485.333) {⑦};
\node[inner sep=0pt,align=center,font=\normalsize,] at (450.000,481.333) {⑧};
\q{104}{769}{241}{12}{三角形：（\hspace{46pt}）}
\q{329}{769}{292}{12}{平行四边形：（\hspace{50pt}）}
\q{626}{769}{200}{12}{梯形：（\hspace{47pt}）}
\hdsection{4}{827}{在下面的点子图中画一画。}
\q{103}{870}{722}{12}{① 画出有一个角是直角的等腰三角形。}
\q{103}{914}{722}{12}{② 画一个梯形，并画出它的一条高，标出上底和下底。}
\q{103}{957}{722}{12}{③ 画一个平行四边形，并画出它的一条高。}
\fill (75.333,673.333) circle (.7bp);
\fill (87.333,673.333) circle (.7bp);
\fill (99.333,673.333) circle (.7bp);
\fill (111.333,673.333) circle (.7bp);
\fill (123.333,673.333) circle (.7bp);
\fill (135.333,673.333) circle (.7bp);
\fill (147.333,673.333) circle (.7bp);
\fill (159.333,673.333) circle (.7bp);
\fill (171.333,673.333) circle (.7bp);
\fill (183.333,673.333) circle (.7bp);
\fill (195.333,673.333) circle (.7bp);
\fill (207.333,673.333) circle (.7bp);
\fill (219.333,673.333) circle (.7bp);
\fill (231.333,673.333) circle (.7bp);
\fill (243.333,673.333) circle (.7bp);
\fill (255.333,673.333) circle (.7bp);
\fill (267.333,673.333) circle (.7bp);
\fill (279.333,673.333) circle (.7bp);
\fill (291.333,673.333) circle (.7bp);
\fill (303.333,673.333) circle (.7bp);
\fill (315.333,673.333) circle (.7bp);
\fill (327.333,673.333) circle (.7bp);
\fill (339.333,673.333) circle (.7bp);
\fill (351.333,673.333) circle (.7bp);
\fill (363.333,673.333) circle (.7bp);
\fill (375.333,673.333) circle (.7bp);
\fill (387.333,673.333) circle (.7bp);
\fill (399.333,673.333) circle (.7bp);
\fill (411.333,673.333) circle (.7bp);
\fill (423.333,673.333) circle (.7bp);
\fill (435.333,673.333) circle (.7bp);
\fill (447.333,673.333) circle (.7bp);
\fill (459.333,673.333) circle (.7bp);
\fill (471.333,673.333) circle (.7bp);
\fill (483.333,673.333) circle (.7bp);
\fill (495.333,673.333) circle (.7bp);
\fill (507.333,673.333) circle (.7bp);
\fill (519.333,673.333) circle (.7bp);
\fill (531.333,673.333) circle (.7bp);
\fill (543.333,673.333) circle (.7bp);
\fill (75.333,685.333) circle (.7bp);
\fill (87.333,685.333) circle (.7bp);
\fill (99.333,685.333) circle (.7bp);
\fill (111.333,685.333) circle (.7bp);
\fill (123.333,685.333) circle (.7bp);
\fill (135.333,685.333) circle (.7bp);
\fill (147.333,685.333) circle (.7bp);
\fill (159.333,685.333) circle (.7bp);
\fill (171.333,685.333) circle (.7bp);
\fill (183.333,685.333) circle (.7bp);
\fill (195.333,685.333) circle (.7bp);
\fill (207.333,685.333) circle (.7bp);
\fill (219.333,685.333) circle (.7bp);
\fill (231.333,685.333) circle (.7bp);
\fill (243.333,685.333) circle (.7bp);
\fill (255.333,685.333) circle (.7bp);
\fill (267.333,685.333) circle (.7bp);
\fill (279.333,685.333) circle (.7bp);
\fill (291.333,685.333) circle (.7bp);
\fill (303.333,685.333) circle (.7bp);
\fill (315.333,685.333) circle (.7bp);
\fill (327.333,685.333) circle (.7bp);
\fill (339.333,685.333) circle (.7bp);
\fill (351.333,685.333) circle (.7bp);
\fill (363.333,685.333) circle (.7bp);
\fill (375.333,685.333) circle (.7bp);
\fill (387.333,685.333) circle (.7bp);
\fill (399.333,685.333) circle (.7bp);
\fill (411.333,685.333) circle (.7bp);
\fill (423.333,685.333) circle (.7bp);
\fill (435.333,685.333) circle (.7bp);
\fill (447.333,685.333) circle (.7bp);
\fill (459.333,685.333) circle (.7bp);
\fill (471.333,685.333) circle (.7bp);
\fill (483.333,685.333) circle (.7bp);
\fill (495.333,685.333) circle (.7bp);
\fill (507.333,685.333) circle (.7bp);
\fill (519.333,685.333) circle (.7bp);
\fill (531.333,685.333) circle (.7bp);
\fill (543.333,685.333) circle (.7bp);
\fill (75.333,697.333) circle (.7bp);
\fill (87.333,697.333) circle (.7bp);
\fill (99.333,697.333) circle (.7bp);
\fill (111.333,697.333) circle (.7bp);
\fill (123.333,697.333) circle (.7bp);
\fill (135.333,697.333) circle (.7bp);
\fill (147.333,697.333) circle (.7bp);
\fill (159.333,697.333) circle (.7bp);
\fill (171.333,697.333) circle (.7bp);
\fill (183.333,697.333) circle (.7bp);
\fill (195.333,697.333) circle (.7bp);
\fill (207.333,697.333) circle (.7bp);
\fill (219.333,697.333) circle (.7bp);
\fill (231.333,697.333) circle (.7bp);
\fill (243.333,697.333) circle (.7bp);
\fill (255.333,697.333) circle (.7bp);
\fill (267.333,697.333) circle (.7bp);
\fill (279.333,697.333) circle (.7bp);
\fill (291.333,697.333) circle (.7bp);
\fill (303.333,697.333) circle (.7bp);
\fill (315.333,697.333) circle (.7bp);
\fill (327.333,697.333) circle (.7bp);
\fill (339.333,697.333) circle (.7bp);
\fill (351.333,697.333) circle (.7bp);
\fill (363.333,697.333) circle (.7bp);
\fill (375.333,697.333) circle (.7bp);
\fill (387.333,697.333) circle (.7bp);
\fill (399.333,697.333) circle (.7bp);
\fill (411.333,697.333) circle (.7bp);
\fill (423.333,697.333) circle (.7bp);
\fill (435.333,697.333) circle (.7bp);
\fill (447.333,697.333) circle (.7bp);
\fill (459.333,697.333) circle (.7bp);
\fill (471.333,697.333) circle (.7bp);
\fill (483.333,697.333) circle (.7bp);
\fill (495.333,697.333) circle (.7bp);
\fill (507.333,697.333) circle (.7bp);
\fill (519.333,697.333) circle (.7bp);
\fill (531.333,697.333) circle (.7bp);
\fill (543.333,697.333) circle (.7bp);
\fill (75.333,709.333) circle (.7bp);
\fill (87.333,709.333) circle (.7bp);
\fill (99.333,709.333) circle (.7bp);
\fill (111.333,709.333) circle (.7bp);
\fill (123.333,709.333) circle (.7bp);
\fill (135.333,709.333) circle (.7bp);
\fill (147.333,709.333) circle (.7bp);
\fill (159.333,709.333) circle (.7bp);
\fill (171.333,709.333) circle (.7bp);
\fill (183.333,709.333) circle (.7bp);
\fill (195.333,709.333) circle (.7bp);
\fill (207.333,709.333) circle (.7bp);
\fill (219.333,709.333) circle (.7bp);
\fill (231.333,709.333) circle (.7bp);
\fill (243.333,709.333) circle (.7bp);
\fill (255.333,709.333) circle (.7bp);
\fill (267.333,709.333) circle (.7bp);
\fill (279.333,709.333) circle (.7bp);
\fill (291.333,709.333) circle (.7bp);
\fill (303.333,709.333) circle (.7bp);
\fill (315.333,709.333) circle (.7bp);
\fill (327.333,709.333) circle (.7bp);
\fill (339.333,709.333) circle (.7bp);
\fill (351.333,709.333) circle (.7bp);
\fill (363.333,709.333) circle (.7bp);
\fill (375.333,709.333) circle (.7bp);
\fill (387.333,709.333) circle (.7bp);
\fill (399.333,709.333) circle (.7bp);
\fill (411.333,709.333) circle (.7bp);
\fill (423.333,709.333) circle (.7bp);
\fill (435.333,709.333) circle (.7bp);
\fill (447.333,709.333) circle (.7bp);
\fill (459.333,709.333) circle (.7bp);
\fill (471.333,709.333) circle (.7bp);
\fill (483.333,709.333) circle (.7bp);
\fill (495.333,709.333) circle (.7bp);
\fill (507.333,709.333) circle (.7bp);
\fill (519.333,709.333) circle (.7bp);
\fill (531.333,709.333) circle (.7bp);
\fill (543.333,709.333) circle (.7bp);
\fill (75.333,721.333) circle (.7bp);
\fill (87.333,721.333) circle (.7bp);
\fill (99.333,721.333) circle (.7bp);
\fill (111.333,721.333) circle (.7bp);
\fill (123.333,721.333) circle (.7bp);
\fill (135.333,721.333) circle (.7bp);
\fill (147.333,721.333) circle (.7bp);
\fill (159.333,721.333) circle (.7bp);
\fill (171.333,721.333) circle (.7bp);
\fill (183.333,721.333) circle (.7bp);
\fill (195.333,721.333) circle (.7bp);
\fill (207.333,721.333) circle (.7bp);
\fill (219.333,721.333) circle (.7bp);
\fill (231.333,721.333) circle (.7bp);
\fill (243.333,721.333) circle (.7bp);
\fill (255.333,721.333) circle (.7bp);
\fill (267.333,721.333) circle (.7bp);
\fill (279.333,721.333) circle (.7bp);
\fill (291.333,721.333) circle (.7bp);
\fill (303.333,721.333) circle (.7bp);
\fill (315.333,721.333) circle (.7bp);
\fill (327.333,721.333) circle (.7bp);
\fill (339.333,721.333) circle (.7bp);
\fill (351.333,721.333) circle (.7bp);
\fill (363.333,721.333) circle (.7bp);
\fill (375.333,721.333) circle (.7bp);
\fill (387.333,721.333) circle (.7bp);
\fill (399.333,721.333) circle (.7bp);
\fill (411.333,721.333) circle (.7bp);
\fill (423.333,721.333) circle (.7bp);
\fill (435.333,721.333) circle (.7bp);
\fill (447.333,721.333) circle (.7bp);
\fill (459.333,721.333) circle (.7bp);
\fill (471.333,721.333) circle (.7bp);
\fill (483.333,721.333) circle (.7bp);
\fill (495.333,721.333) circle (.7bp);
\fill (507.333,721.333) circle (.7bp);
\fill (519.333,721.333) circle (.7bp);
\fill (531.333,721.333) circle (.7bp);
\fill (543.333,721.333) circle (.7bp);
\fill (75.333,733.333) circle (.7bp);
\fill (87.333,733.333) circle (.7bp);
\fill (99.333,733.333) circle (.7bp);
\fill (111.333,733.333) circle (.7bp);
\fill (123.333,733.333) circle (.7bp);
\fill (135.333,733.333) circle (.7bp);
\fill (147.333,733.333) circle (.7bp);
\fill (159.333,733.333) circle (.7bp);
\fill (171.333,733.333) circle (.7bp);
\fill (183.333,733.333) circle (.7bp);
\fill (195.333,733.333) circle (.7bp);
\fill (207.333,733.333) circle (.7bp);
\fill (219.333,733.333) circle (.7bp);
\fill (231.333,733.333) circle (.7bp);
\fill (243.333,733.333) circle (.7bp);
\fill (255.333,733.333) circle (.7bp);
\fill (267.333,733.333) circle (.7bp);
\fill (279.333,733.333) circle (.7bp);
\fill (291.333,733.333) circle (.7bp);
\fill (303.333,733.333) circle (.7bp);
\fill (315.333,733.333) circle (.7bp);
\fill (327.333,733.333) circle (.7bp);
\fill (339.333,733.333) circle (.7bp);
\fill (351.333,733.333) circle (.7bp);
\fill (363.333,733.333) circle (.7bp);
\fill (375.333,733.333) circle (.7bp);
\fill (387.333,733.333) circle (.7bp);
\fill (399.333,733.333) circle (.7bp);
\fill (411.333,733.333) circle (.7bp);
\fill (423.333,733.333) circle (.7bp);
\fill (435.333,733.333) circle (.7bp);
\fill (447.333,733.333) circle (.7bp);
\fill (459.333,733.333) circle (.7bp);
\fill (471.333,733.333) circle (.7bp);
\fill (483.333,733.333) circle (.7bp);
\fill (495.333,733.333) circle (.7bp);
\fill (507.333,733.333) circle (.7bp);
\fill (519.333,733.333) circle (.7bp);
\fill (531.333,733.333) circle (.7bp);
\fill (543.333,733.333) circle (.7bp);
\fill (75.333,745.333) circle (.7bp);
\fill (87.333,745.333) circle (.7bp);
\fill (99.333,745.333) circle (.7bp);
\fill (111.333,745.333) circle (.7bp);
\fill (123.333,745.333) circle (.7bp);
\fill (135.333,745.333) circle (.7bp);
\fill (147.333,745.333) circle (.7bp);
\fill (159.333,745.333) circle (.7bp);
\fill (171.333,745.333) circle (.7bp);
\fill (183.333,745.333) circle (.7bp);
\fill (195.333,745.333) circle (.7bp);
\fill (207.333,745.333) circle (.7bp);
\fill (219.333,745.333) circle (.7bp);
\fill (231.333,745.333) circle (.7bp);
\fill (243.333,745.333) circle (.7bp);
\fill (255.333,745.333) circle (.7bp);
\fill (267.333,745.333) circle (.7bp);
\fill (279.333,745.333) circle (.7bp);
\fill (291.333,745.333) circle (.7bp);
\fill (303.333,745.333) circle (.7bp);
\fill (315.333,745.333) circle (.7bp);
\fill (327.333,745.333) circle (.7bp);
\fill (339.333,745.333) circle (.7bp);
\fill (351.333,745.333) circle (.7bp);
\fill (363.333,745.333) circle (.7bp);
\fill (375.333,745.333) circle (.7bp);
\fill (387.333,745.333) circle (.7bp);
\fill (399.333,745.333) circle (.7bp);
\fill (411.333,745.333) circle (.7bp);
\fill (423.333,745.333) circle (.7bp);
\fill (435.333,745.333) circle (.7bp);
\fill (447.333,745.333) circle (.7bp);
\fill (459.333,745.333) circle (.7bp);
\fill (471.333,745.333) circle (.7bp);
\fill (483.333,745.333) circle (.7bp);
\fill (495.333,745.333) circle (.7bp);
\fill (507.333,745.333) circle (.7bp);
\fill (519.333,745.333) circle (.7bp);
\fill (531.333,745.333) circle (.7bp);
\fill (543.333,745.333) circle (.7bp);
\fill (75.333,757.333) circle (.7bp);
\fill (87.333,757.333) circle (.7bp);
\fill (99.333,757.333) circle (.7bp);
\fill (111.333,757.333) circle (.7bp);
\fill (123.333,757.333) circle (.7bp);
\fill (135.333,757.333) circle (.7bp);
\fill (147.333,757.333) circle (.7bp);
\fill (159.333,757.333) circle (.7bp);
\fill (171.333,757.333) circle (.7bp);
\fill (183.333,757.333) circle (.7bp);
\fill (195.333,757.333) circle (.7bp);
\fill (207.333,757.333) circle (.7bp);
\fill (219.333,757.333) circle (.7bp);
\fill (231.333,757.333) circle (.7bp);
\fill (243.333,757.333) circle (.7bp);
\fill (255.333,757.333) circle (.7bp);
\fill (267.333,757.333) circle (.7bp);
\fill (279.333,757.333) circle (.7bp);
\fill (291.333,757.333) circle (.7bp);
\fill (303.333,757.333) circle (.7bp);
\fill (315.333,757.333) circle (.7bp);
\fill (327.333,757.333) circle (.7bp);
\fill (339.333,757.333) circle (.7bp);
\fill (351.333,757.333) circle (.7bp);
\fill (363.333,757.333) circle (.7bp);
\fill (375.333,757.333) circle (.7bp);
\fill (387.333,757.333) circle (.7bp);
\fill (399.333,757.333) circle (.7bp);
\fill (411.333,757.333) circle (.7bp);
\fill (423.333,757.333) circle (.7bp);
\fill (435.333,757.333) circle (.7bp);
\fill (447.333,757.333) circle (.7bp);
\fill (459.333,757.333) circle (.7bp);
\fill (471.333,757.333) circle (.7bp);
\fill (483.333,757.333) circle (.7bp);
\fill (495.333,757.333) circle (.7bp);
\fill (507.333,757.333) circle (.7bp);
\fill (519.333,757.333) circle (.7bp);
\fill (531.333,757.333) circle (.7bp);
\fill (543.333,757.333) circle (.7bp);
\fill (75.333,769.333) circle (.7bp);
\fill (87.333,769.333) circle (.7bp);
\fill (99.333,769.333) circle (.7bp);
\fill (111.333,769.333) circle (.7bp);
\fill (123.333,769.333) circle (.7bp);
\fill (135.333,769.333) circle (.7bp);
\fill (147.333,769.333) circle (.7bp);
\fill (159.333,769.333) circle (.7bp);
\fill (171.333,769.333) circle (.7bp);
\fill (183.333,769.333) circle (.7bp);
\fill (195.333,769.333) circle (.7bp);
\fill (207.333,769.333) circle (.7bp);
\fill (219.333,769.333) circle (.7bp);
\fill (231.333,769.333) circle (.7bp);
\fill (243.333,769.333) circle (.7bp);
\fill (255.333,769.333) circle (.7bp);
\fill (267.333,769.333) circle (.7bp);
\fill (279.333,769.333) circle (.7bp);
\fill (291.333,769.333) circle (.7bp);
\fill (303.333,769.333) circle (.7bp);
\fill (315.333,769.333) circle (.7bp);
\fill (327.333,769.333) circle (.7bp);
\fill (339.333,769.333) circle (.7bp);
\fill (351.333,769.333) circle (.7bp);
\fill (363.333,769.333) circle (.7bp);
\fill (375.333,769.333) circle (.7bp);
\fill (387.333,769.333) circle (.7bp);
\fill (399.333,769.333) circle (.7bp);
\fill (411.333,769.333) circle (.7bp);
\fill (423.333,769.333) circle (.7bp);
\fill (435.333,769.333) circle (.7bp);
\fill (447.333,769.333) circle (.7bp);
\fill (459.333,769.333) circle (.7bp);
\fill (471.333,769.333) circle (.7bp);
\fill (483.333,769.333) circle (.7bp);
\fill (495.333,769.333) circle (.7bp);
\fill (507.333,769.333) circle (.7bp);
\fill (519.333,769.333) circle (.7bp);
\fill (531.333,769.333) circle (.7bp);
\fill (543.333,769.333) circle (.7bp);

\begin{scope}[color=answerred,draw=answerred,text=answerred]% 高清文字版专用填空位置；原页答案另存于 ../overlays。
\ans{268.667}{186.667}{40.000}{14}{$\checkmark$}
\ans{352.000}{186.667}{40.000}{14}{$\checkmark$}
\ans{162.000}{263.333}{32.000}{12}{$5$}
\ans{162.000}{292.667}{32.000}{12}{$3$}
\ans{420.667}{263.333}{32.000}{12}{$3$}
\ans{420.667}{292.667}{32.000}{12}{$1$}
\ans{129.613}{512.343}{66.000}{11}{②③⑤}
\ans{305.517}{512.343}{70.000}{11}{①⑥⑧}
\ans{477.115}{512.343}{43.333}{11}{④}
\begin{scope}[scale=0.6666667]\draw[answerred,line width=1pt] (149,1136)--(149,1064)--(221,1136)--cycle;\draw (149,1125)--(160,1125)--(160,1136);\draw (144,1100)--(154,1100) (185,1131)--(185,1141);\node[answerred,above] at (185,1050) {①};\draw[answerred,line width=1pt] (365,1136)--(401,1064)--(473,1064)--(509,1136)--cycle;\draw[answerred,dashed] (401,1064)--(401,1136);\draw (401,1125)--(412,1125)--(412,1136);\node[above] at (437,1061) {上底};\node[below] at (437,1136) {下底};\node[left] at (401,1100) {高};\node[above] at (365,1050) {②};\draw[answerred,line width=1pt] (617,1136)--(653,1064)--(761,1064)--(725,1136)--cycle;\draw[answerred,dashed] (653,1064)--(653,1136);\draw (653,1125)--(664,1125)--(664,1136);\node[left] at (653,1100) {高};\node[above] at (616,1050) {③};\end{scope}
\end{scope}
\end{scope}\pgfresetboundingbox\path[use as bounding box] (56.0,248.0) rectangle (527.0,324.67);
\end{tikzpicture}}\par



\par

\par\noindent

\QHeading{第2题　数图形的不同办法}

\Answer{方法一：按大小或包含的小块数量分类数，再合计。方法二：固定一个顶点或一条底边，逐一寻找另外的边，并把每个已经计入的图形编号。无论采用哪种办法，都要把由多个小图形拼成的大图形计入，同时避免把同一个图形反复计数。}

\QHeading{第3题　分一分。（填序号。）}

\par\noindent\makebox[\linewidth][c]{
\begin{tikzpicture}[x=1bp,y=-1bp,color=black,draw=black,text=black]
\begin{scope}\clip (56.0,358.0) rectangle (516.0,549.33);
\hdhead{94}{总复习}
\hdsection{1}{150}{在既有平行线段又有垂直线段的图形下面的（\quad）里画“$\checkmark$”。}
\draw[line width=.7pt] (148.000,168.000)--(170.000,136.667)--(220.000,136.667)--(197.333,168.000)--cycle;
\draw[line width=.7pt] (242.667,136.667)--(308.667,136.667)--(308.667,168.000)--(242.667,168.000)--cycle;
\draw[line width=.7pt] (334.667,136.667)--(372.000,136.667)--(385.333,168.000)--(334.667,168.000)--cycle;
\draw[line width=.7pt] (406.667,136.667)--(406.667,168.000)--(462.667,168.000)--cycle;
\q{236}{279}{110}{12}{（\hspace{30pt}）}
\q{374}{279}{110}{12}{（\hspace{30pt}）}
\q{515}{279}{110}{12}{（\hspace{30pt}）}
\q{615}{279}{110}{12}{（\hspace{30pt}）}
\hdsection{2}{335}{数一数，填一填。}
\q{102}{380}{42}{13}{①}
\q{403}{380}{42}{13}{②}
\draw[line width=.7pt] (89.333,302.000)--(116.667,256.000)--(143.333,302.000)--cycle;
\draw[line width=.7pt] (103.000,279.000)--(130.000,279.000)--(116.667,302.000)--cycle;
\draw[line width=.7pt] (290.000,302.000)--(316.667,256.000)--(374.667,256.000)--(402.000,302.000)--cycle;
\draw[line width=.7pt] (316.667,256.000)--(343.333,302.000);
\draw[line width=.7pt] (303.333,279.000)--(330.000,279.000)--(330.000,302.000);
\q{215}{387}{196}{12}{（\hspace{25pt}）个三角形}
\q{215}{431}{196}{12}{（\hspace{25pt}）个梯形}
\q{603}{387}{230}{12}{（\hspace{25pt}）个梯形}
\q{603}{431}{250}{12}{（\hspace{25pt}）个平行四边形}
\hdsection{3}{490}{分一分。（填序号。）}
\draw[line width=.7pt] (139.333,396.000)--(152.667,363.333)--(191.333,363.333)--(177.333,396.000)--cycle;
\draw[line width=.7pt] (245.333,395.333)--(260.667,364.667)--(275.333,395.333)--cycle;
\draw[line width=.7pt] (326.000,363.333)--(326.000,396.000)--(384.667,396.000)--cycle;
\draw[line width=.7pt] (416.667,394.000)--(432.667,364.667)--(466.667,364.667)--(482.667,394.000)--cycle;
\draw[line width=.7pt] (146.000,468.000)--(164.667,434.000)--(184.667,468.000)--cycle;
\draw[line width=.7pt] (238.667,434.667)--(282.667,434.667)--(282.667,468.000)--(238.667,468.000)--cycle;
\draw[line width=.7pt] (338.000,448.000)--(354.667,434.667)--(372.000,448.000)--(366.667,468.000)--(344.000,468.000)--cycle;
\draw[line width=.7pt] (439.333,467.333)--(444.000,434.000)--(460.667,434.000)--(456.000,467.333)--cycle;
\node[inner sep=0pt,align=center,font=\normalsize,] at (165.333,414.000) {①};
\node[inner sep=0pt,align=center,font=\normalsize,] at (260.667,408.667) {②};
\node[inner sep=0pt,align=center,font=\normalsize,] at (355.333,413.333) {③};
\node[inner sep=0pt,align=center,font=\normalsize,] at (450.000,408.667) {④};
\node[inner sep=0pt,align=center,font=\normalsize,] at (165.333,485.333) {⑤};
\node[inner sep=0pt,align=center,font=\normalsize,] at (260.667,485.333) {⑥};
\node[inner sep=0pt,align=center,font=\normalsize,] at (355.333,485.333) {⑦};
\node[inner sep=0pt,align=center,font=\normalsize,] at (450.000,481.333) {⑧};
\q{104}{769}{241}{12}{三角形：（\hspace{46pt}）}
\q{329}{769}{292}{12}{平行四边形：（\hspace{50pt}）}
\q{626}{769}{200}{12}{梯形：（\hspace{47pt}）}
\hdsection{4}{827}{在下面的点子图中画一画。}
\q{103}{870}{722}{12}{① 画出有一个角是直角的等腰三角形。}
\q{103}{914}{722}{12}{② 画一个梯形，并画出它的一条高，标出上底和下底。}
\q{103}{957}{722}{12}{③ 画一个平行四边形，并画出它的一条高。}
\fill (75.333,673.333) circle (.7bp);
\fill (87.333,673.333) circle (.7bp);
\fill (99.333,673.333) circle (.7bp);
\fill (111.333,673.333) circle (.7bp);
\fill (123.333,673.333) circle (.7bp);
\fill (135.333,673.333) circle (.7bp);
\fill (147.333,673.333) circle (.7bp);
\fill (159.333,673.333) circle (.7bp);
\fill (171.333,673.333) circle (.7bp);
\fill (183.333,673.333) circle (.7bp);
\fill (195.333,673.333) circle (.7bp);
\fill (207.333,673.333) circle (.7bp);
\fill (219.333,673.333) circle (.7bp);
\fill (231.333,673.333) circle (.7bp);
\fill (243.333,673.333) circle (.7bp);
\fill (255.333,673.333) circle (.7bp);
\fill (267.333,673.333) circle (.7bp);
\fill (279.333,673.333) circle (.7bp);
\fill (291.333,673.333) circle (.7bp);
\fill (303.333,673.333) circle (.7bp);
\fill (315.333,673.333) circle (.7bp);
\fill (327.333,673.333) circle (.7bp);
\fill (339.333,673.333) circle (.7bp);
\fill (351.333,673.333) circle (.7bp);
\fill (363.333,673.333) circle (.7bp);
\fill (375.333,673.333) circle (.7bp);
\fill (387.333,673.333) circle (.7bp);
\fill (399.333,673.333) circle (.7bp);
\fill (411.333,673.333) circle (.7bp);
\fill (423.333,673.333) circle (.7bp);
\fill (435.333,673.333) circle (.7bp);
\fill (447.333,673.333) circle (.7bp);
\fill (459.333,673.333) circle (.7bp);
\fill (471.333,673.333) circle (.7bp);
\fill (483.333,673.333) circle (.7bp);
\fill (495.333,673.333) circle (.7bp);
\fill (507.333,673.333) circle (.7bp);
\fill (519.333,673.333) circle (.7bp);
\fill (531.333,673.333) circle (.7bp);
\fill (543.333,673.333) circle (.7bp);
\fill (75.333,685.333) circle (.7bp);
\fill (87.333,685.333) circle (.7bp);
\fill (99.333,685.333) circle (.7bp);
\fill (111.333,685.333) circle (.7bp);
\fill (123.333,685.333) circle (.7bp);
\fill (135.333,685.333) circle (.7bp);
\fill (147.333,685.333) circle (.7bp);
\fill (159.333,685.333) circle (.7bp);
\fill (171.333,685.333) circle (.7bp);
\fill (183.333,685.333) circle (.7bp);
\fill (195.333,685.333) circle (.7bp);
\fill (207.333,685.333) circle (.7bp);
\fill (219.333,685.333) circle (.7bp);
\fill (231.333,685.333) circle (.7bp);
\fill (243.333,685.333) circle (.7bp);
\fill (255.333,685.333) circle (.7bp);
\fill (267.333,685.333) circle (.7bp);
\fill (279.333,685.333) circle (.7bp);
\fill (291.333,685.333) circle (.7bp);
\fill (303.333,685.333) circle (.7bp);
\fill (315.333,685.333) circle (.7bp);
\fill (327.333,685.333) circle (.7bp);
\fill (339.333,685.333) circle (.7bp);
\fill (351.333,685.333) circle (.7bp);
\fill (363.333,685.333) circle (.7bp);
\fill (375.333,685.333) circle (.7bp);
\fill (387.333,685.333) circle (.7bp);
\fill (399.333,685.333) circle (.7bp);
\fill (411.333,685.333) circle (.7bp);
\fill (423.333,685.333) circle (.7bp);
\fill (435.333,685.333) circle (.7bp);
\fill (447.333,685.333) circle (.7bp);
\fill (459.333,685.333) circle (.7bp);
\fill (471.333,685.333) circle (.7bp);
\fill (483.333,685.333) circle (.7bp);
\fill (495.333,685.333) circle (.7bp);
\fill (507.333,685.333) circle (.7bp);
\fill (519.333,685.333) circle (.7bp);
\fill (531.333,685.333) circle (.7bp);
\fill (543.333,685.333) circle (.7bp);
\fill (75.333,697.333) circle (.7bp);
\fill (87.333,697.333) circle (.7bp);
\fill (99.333,697.333) circle (.7bp);
\fill (111.333,697.333) circle (.7bp);
\fill (123.333,697.333) circle (.7bp);
\fill (135.333,697.333) circle (.7bp);
\fill (147.333,697.333) circle (.7bp);
\fill (159.333,697.333) circle (.7bp);
\fill (171.333,697.333) circle (.7bp);
\fill (183.333,697.333) circle (.7bp);
\fill (195.333,697.333) circle (.7bp);
\fill (207.333,697.333) circle (.7bp);
\fill (219.333,697.333) circle (.7bp);
\fill (231.333,697.333) circle (.7bp);
\fill (243.333,697.333) circle (.7bp);
\fill (255.333,697.333) circle (.7bp);
\fill (267.333,697.333) circle (.7bp);
\fill (279.333,697.333) circle (.7bp);
\fill (291.333,697.333) circle (.7bp);
\fill (303.333,697.333) circle (.7bp);
\fill (315.333,697.333) circle (.7bp);
\fill (327.333,697.333) circle (.7bp);
\fill (339.333,697.333) circle (.7bp);
\fill (351.333,697.333) circle (.7bp);
\fill (363.333,697.333) circle (.7bp);
\fill (375.333,697.333) circle (.7bp);
\fill (387.333,697.333) circle (.7bp);
\fill (399.333,697.333) circle (.7bp);
\fill (411.333,697.333) circle (.7bp);
\fill (423.333,697.333) circle (.7bp);
\fill (435.333,697.333) circle (.7bp);
\fill (447.333,697.333) circle (.7bp);
\fill (459.333,697.333) circle (.7bp);
\fill (471.333,697.333) circle (.7bp);
\fill (483.333,697.333) circle (.7bp);
\fill (495.333,697.333) circle (.7bp);
\fill (507.333,697.333) circle (.7bp);
\fill (519.333,697.333) circle (.7bp);
\fill (531.333,697.333) circle (.7bp);
\fill (543.333,697.333) circle (.7bp);
\fill (75.333,709.333) circle (.7bp);
\fill (87.333,709.333) circle (.7bp);
\fill (99.333,709.333) circle (.7bp);
\fill (111.333,709.333) circle (.7bp);
\fill (123.333,709.333) circle (.7bp);
\fill (135.333,709.333) circle (.7bp);
\fill (147.333,709.333) circle (.7bp);
\fill (159.333,709.333) circle (.7bp);
\fill (171.333,709.333) circle (.7bp);
\fill (183.333,709.333) circle (.7bp);
\fill (195.333,709.333) circle (.7bp);
\fill (207.333,709.333) circle (.7bp);
\fill (219.333,709.333) circle (.7bp);
\fill (231.333,709.333) circle (.7bp);
\fill (243.333,709.333) circle (.7bp);
\fill (255.333,709.333) circle (.7bp);
\fill (267.333,709.333) circle (.7bp);
\fill (279.333,709.333) circle (.7bp);
\fill (291.333,709.333) circle (.7bp);
\fill (303.333,709.333) circle (.7bp);
\fill (315.333,709.333) circle (.7bp);
\fill (327.333,709.333) circle (.7bp);
\fill (339.333,709.333) circle (.7bp);
\fill (351.333,709.333) circle (.7bp);
\fill (363.333,709.333) circle (.7bp);
\fill (375.333,709.333) circle (.7bp);
\fill (387.333,709.333) circle (.7bp);
\fill (399.333,709.333) circle (.7bp);
\fill (411.333,709.333) circle (.7bp);
\fill (423.333,709.333) circle (.7bp);
\fill (435.333,709.333) circle (.7bp);
\fill (447.333,709.333) circle (.7bp);
\fill (459.333,709.333) circle (.7bp);
\fill (471.333,709.333) circle (.7bp);
\fill (483.333,709.333) circle (.7bp);
\fill (495.333,709.333) circle (.7bp);
\fill (507.333,709.333) circle (.7bp);
\fill (519.333,709.333) circle (.7bp);
\fill (531.333,709.333) circle (.7bp);
\fill (543.333,709.333) circle (.7bp);
\fill (75.333,721.333) circle (.7bp);
\fill (87.333,721.333) circle (.7bp);
\fill (99.333,721.333) circle (.7bp);
\fill (111.333,721.333) circle (.7bp);
\fill (123.333,721.333) circle (.7bp);
\fill (135.333,721.333) circle (.7bp);
\fill (147.333,721.333) circle (.7bp);
\fill (159.333,721.333) circle (.7bp);
\fill (171.333,721.333) circle (.7bp);
\fill (183.333,721.333) circle (.7bp);
\fill (195.333,721.333) circle (.7bp);
\fill (207.333,721.333) circle (.7bp);
\fill (219.333,721.333) circle (.7bp);
\fill (231.333,721.333) circle (.7bp);
\fill (243.333,721.333) circle (.7bp);
\fill (255.333,721.333) circle (.7bp);
\fill (267.333,721.333) circle (.7bp);
\fill (279.333,721.333) circle (.7bp);
\fill (291.333,721.333) circle (.7bp);
\fill (303.333,721.333) circle (.7bp);
\fill (315.333,721.333) circle (.7bp);
\fill (327.333,721.333) circle (.7bp);
\fill (339.333,721.333) circle (.7bp);
\fill (351.333,721.333) circle (.7bp);
\fill (363.333,721.333) circle (.7bp);
\fill (375.333,721.333) circle (.7bp);
\fill (387.333,721.333) circle (.7bp);
\fill (399.333,721.333) circle (.7bp);
\fill (411.333,721.333) circle (.7bp);
\fill (423.333,721.333) circle (.7bp);
\fill (435.333,721.333) circle (.7bp);
\fill (447.333,721.333) circle (.7bp);
\fill (459.333,721.333) circle (.7bp);
\fill (471.333,721.333) circle (.7bp);
\fill (483.333,721.333) circle (.7bp);
\fill (495.333,721.333) circle (.7bp);
\fill (507.333,721.333) circle (.7bp);
\fill (519.333,721.333) circle (.7bp);
\fill (531.333,721.333) circle (.7bp);
\fill (543.333,721.333) circle (.7bp);
\fill (75.333,733.333) circle (.7bp);
\fill (87.333,733.333) circle (.7bp);
\fill (99.333,733.333) circle (.7bp);
\fill (111.333,733.333) circle (.7bp);
\fill (123.333,733.333) circle (.7bp);
\fill (135.333,733.333) circle (.7bp);
\fill (147.333,733.333) circle (.7bp);
\fill (159.333,733.333) circle (.7bp);
\fill (171.333,733.333) circle (.7bp);
\fill (183.333,733.333) circle (.7bp);
\fill (195.333,733.333) circle (.7bp);
\fill (207.333,733.333) circle (.7bp);
\fill (219.333,733.333) circle (.7bp);
\fill (231.333,733.333) circle (.7bp);
\fill (243.333,733.333) circle (.7bp);
\fill (255.333,733.333) circle (.7bp);
\fill (267.333,733.333) circle (.7bp);
\fill (279.333,733.333) circle (.7bp);
\fill (291.333,733.333) circle (.7bp);
\fill (303.333,733.333) circle (.7bp);
\fill (315.333,733.333) circle (.7bp);
\fill (327.333,733.333) circle (.7bp);
\fill (339.333,733.333) circle (.7bp);
\fill (351.333,733.333) circle (.7bp);
\fill (363.333,733.333) circle (.7bp);
\fill (375.333,733.333) circle (.7bp);
\fill (387.333,733.333) circle (.7bp);
\fill (399.333,733.333) circle (.7bp);
\fill (411.333,733.333) circle (.7bp);
\fill (423.333,733.333) circle (.7bp);
\fill (435.333,733.333) circle (.7bp);
\fill (447.333,733.333) circle (.7bp);
\fill (459.333,733.333) circle (.7bp);
\fill (471.333,733.333) circle (.7bp);
\fill (483.333,733.333) circle (.7bp);
\fill (495.333,733.333) circle (.7bp);
\fill (507.333,733.333) circle (.7bp);
\fill (519.333,733.333) circle (.7bp);
\fill (531.333,733.333) circle (.7bp);
\fill (543.333,733.333) circle (.7bp);
\fill (75.333,745.333) circle (.7bp);
\fill (87.333,745.333) circle (.7bp);
\fill (99.333,745.333) circle (.7bp);
\fill (111.333,745.333) circle (.7bp);
\fill (123.333,745.333) circle (.7bp);
\fill (135.333,745.333) circle (.7bp);
\fill (147.333,745.333) circle (.7bp);
\fill (159.333,745.333) circle (.7bp);
\fill (171.333,745.333) circle (.7bp);
\fill (183.333,745.333) circle (.7bp);
\fill (195.333,745.333) circle (.7bp);
\fill (207.333,745.333) circle (.7bp);
\fill (219.333,745.333) circle (.7bp);
\fill (231.333,745.333) circle (.7bp);
\fill (243.333,745.333) circle (.7bp);
\fill (255.333,745.333) circle (.7bp);
\fill (267.333,745.333) circle (.7bp);
\fill (279.333,745.333) circle (.7bp);
\fill (291.333,745.333) circle (.7bp);
\fill (303.333,745.333) circle (.7bp);
\fill (315.333,745.333) circle (.7bp);
\fill (327.333,745.333) circle (.7bp);
\fill (339.333,745.333) circle (.7bp);
\fill (351.333,745.333) circle (.7bp);
\fill (363.333,745.333) circle (.7bp);
\fill (375.333,745.333) circle (.7bp);
\fill (387.333,745.333) circle (.7bp);
\fill (399.333,745.333) circle (.7bp);
\fill (411.333,745.333) circle (.7bp);
\fill (423.333,745.333) circle (.7bp);
\fill (435.333,745.333) circle (.7bp);
\fill (447.333,745.333) circle (.7bp);
\fill (459.333,745.333) circle (.7bp);
\fill (471.333,745.333) circle (.7bp);
\fill (483.333,745.333) circle (.7bp);
\fill (495.333,745.333) circle (.7bp);
\fill (507.333,745.333) circle (.7bp);
\fill (519.333,745.333) circle (.7bp);
\fill (531.333,745.333) circle (.7bp);
\fill (543.333,745.333) circle (.7bp);
\fill (75.333,757.333) circle (.7bp);
\fill (87.333,757.333) circle (.7bp);
\fill (99.333,757.333) circle (.7bp);
\fill (111.333,757.333) circle (.7bp);
\fill (123.333,757.333) circle (.7bp);
\fill (135.333,757.333) circle (.7bp);
\fill (147.333,757.333) circle (.7bp);
\fill (159.333,757.333) circle (.7bp);
\fill (171.333,757.333) circle (.7bp);
\fill (183.333,757.333) circle (.7bp);
\fill (195.333,757.333) circle (.7bp);
\fill (207.333,757.333) circle (.7bp);
\fill (219.333,757.333) circle (.7bp);
\fill (231.333,757.333) circle (.7bp);
\fill (243.333,757.333) circle (.7bp);
\fill (255.333,757.333) circle (.7bp);
\fill (267.333,757.333) circle (.7bp);
\fill (279.333,757.333) circle (.7bp);
\fill (291.333,757.333) circle (.7bp);
\fill (303.333,757.333) circle (.7bp);
\fill (315.333,757.333) circle (.7bp);
\fill (327.333,757.333) circle (.7bp);
\fill (339.333,757.333) circle (.7bp);
\fill (351.333,757.333) circle (.7bp);
\fill (363.333,757.333) circle (.7bp);
\fill (375.333,757.333) circle (.7bp);
\fill (387.333,757.333) circle (.7bp);
\fill (399.333,757.333) circle (.7bp);
\fill (411.333,757.333) circle (.7bp);
\fill (423.333,757.333) circle (.7bp);
\fill (435.333,757.333) circle (.7bp);
\fill (447.333,757.333) circle (.7bp);
\fill (459.333,757.333) circle (.7bp);
\fill (471.333,757.333) circle (.7bp);
\fill (483.333,757.333) circle (.7bp);
\fill (495.333,757.333) circle (.7bp);
\fill (507.333,757.333) circle (.7bp);
\fill (519.333,757.333) circle (.7bp);
\fill (531.333,757.333) circle (.7bp);
\fill (543.333,757.333) circle (.7bp);
\fill (75.333,769.333) circle (.7bp);
\fill (87.333,769.333) circle (.7bp);
\fill (99.333,769.333) circle (.7bp);
\fill (111.333,769.333) circle (.7bp);
\fill (123.333,769.333) circle (.7bp);
\fill (135.333,769.333) circle (.7bp);
\fill (147.333,769.333) circle (.7bp);
\fill (159.333,769.333) circle (.7bp);
\fill (171.333,769.333) circle (.7bp);
\fill (183.333,769.333) circle (.7bp);
\fill (195.333,769.333) circle (.7bp);
\fill (207.333,769.333) circle (.7bp);
\fill (219.333,769.333) circle (.7bp);
\fill (231.333,769.333) circle (.7bp);
\fill (243.333,769.333) circle (.7bp);
\fill (255.333,769.333) circle (.7bp);
\fill (267.333,769.333) circle (.7bp);
\fill (279.333,769.333) circle (.7bp);
\fill (291.333,769.333) circle (.7bp);
\fill (303.333,769.333) circle (.7bp);
\fill (315.333,769.333) circle (.7bp);
\fill (327.333,769.333) circle (.7bp);
\fill (339.333,769.333) circle (.7bp);
\fill (351.333,769.333) circle (.7bp);
\fill (363.333,769.333) circle (.7bp);
\fill (375.333,769.333) circle (.7bp);
\fill (387.333,769.333) circle (.7bp);
\fill (399.333,769.333) circle (.7bp);
\fill (411.333,769.333) circle (.7bp);
\fill (423.333,769.333) circle (.7bp);
\fill (435.333,769.333) circle (.7bp);
\fill (447.333,769.333) circle (.7bp);
\fill (459.333,769.333) circle (.7bp);
\fill (471.333,769.333) circle (.7bp);
\fill (483.333,769.333) circle (.7bp);
\fill (495.333,769.333) circle (.7bp);
\fill (507.333,769.333) circle (.7bp);
\fill (519.333,769.333) circle (.7bp);
\fill (531.333,769.333) circle (.7bp);
\fill (543.333,769.333) circle (.7bp);

\begin{scope}[color=answerred,draw=answerred,text=answerred]% 高清文字版专用填空位置；原页答案另存于 ../overlays。
\ans{268.667}{186.667}{40.000}{14}{$\checkmark$}
\ans{352.000}{186.667}{40.000}{14}{$\checkmark$}
\ans{162.000}{263.333}{32.000}{12}{$5$}
\ans{162.000}{292.667}{32.000}{12}{$3$}
\ans{420.667}{263.333}{32.000}{12}{$3$}
\ans{420.667}{292.667}{32.000}{12}{$1$}
\ans{129.613}{512.343}{66.000}{11}{②③⑤}
\ans{305.517}{512.343}{70.000}{11}{①⑥⑧}
\ans{477.115}{512.343}{43.333}{11}{④}
\begin{scope}[scale=0.6666667]\draw[answerred,line width=1pt] (149,1136)--(149,1064)--(221,1136)--cycle;\draw (149,1125)--(160,1125)--(160,1136);\draw (144,1100)--(154,1100) (185,1131)--(185,1141);\node[answerred,above] at (185,1050) {①};\draw[answerred,line width=1pt] (365,1136)--(401,1064)--(473,1064)--(509,1136)--cycle;\draw[answerred,dashed] (401,1064)--(401,1136);\draw (401,1125)--(412,1125)--(412,1136);\node[above] at (437,1061) {上底};\node[below] at (437,1136) {下底};\node[left] at (401,1100) {高};\node[above] at (365,1050) {②};\draw[answerred,line width=1pt] (617,1136)--(653,1064)--(761,1064)--(725,1136)--cycle;\draw[answerred,dashed] (653,1064)--(653,1136);\draw (653,1125)--(664,1125)--(664,1136);\node[left] at (653,1100) {高};\node[above] at (616,1050) {③};\end{scope}
\end{scope}
\end{scope}\pgfresetboundingbox\path[use as bounding box] (56.0,358.0) rectangle (516.0,549.33);
\end{tikzpicture}}\par



\par



\par\noindent

\QHeading{第4题　在下面的点子图中画一画。}

\par\noindent\makebox[\linewidth][c]{
\begin{tikzpicture}[x=1bp,y=-1bp,color=black,draw=black,text=black]
\begin{scope}\clip (69.0,667.0) rectangle (549.0,776.0);
\hdhead{94}{总复习}
\hdsection{1}{150}{在既有平行线段又有垂直线段的图形下面的（\quad）里画“$\checkmark$”。}
\draw[line width=.7pt] (148.000,168.000)--(170.000,136.667)--(220.000,136.667)--(197.333,168.000)--cycle;
\draw[line width=.7pt] (242.667,136.667)--(308.667,136.667)--(308.667,168.000)--(242.667,168.000)--cycle;
\draw[line width=.7pt] (334.667,136.667)--(372.000,136.667)--(385.333,168.000)--(334.667,168.000)--cycle;
\draw[line width=.7pt] (406.667,136.667)--(406.667,168.000)--(462.667,168.000)--cycle;
\q{236}{279}{110}{12}{（\hspace{30pt}）}
\q{374}{279}{110}{12}{（\hspace{30pt}）}
\q{515}{279}{110}{12}{（\hspace{30pt}）}
\q{615}{279}{110}{12}{（\hspace{30pt}）}
\hdsection{2}{335}{数一数，填一填。}
\q{102}{380}{42}{13}{①}
\q{403}{380}{42}{13}{②}
\draw[line width=.7pt] (89.333,302.000)--(116.667,256.000)--(143.333,302.000)--cycle;
\draw[line width=.7pt] (103.000,279.000)--(130.000,279.000)--(116.667,302.000)--cycle;
\draw[line width=.7pt] (290.000,302.000)--(316.667,256.000)--(374.667,256.000)--(402.000,302.000)--cycle;
\draw[line width=.7pt] (316.667,256.000)--(343.333,302.000);
\draw[line width=.7pt] (303.333,279.000)--(330.000,279.000)--(330.000,302.000);
\q{215}{387}{196}{12}{（\hspace{25pt}）个三角形}
\q{215}{431}{196}{12}{（\hspace{25pt}）个梯形}
\q{603}{387}{230}{12}{（\hspace{25pt}）个梯形}
\q{603}{431}{250}{12}{（\hspace{25pt}）个平行四边形}
\hdsection{3}{490}{分一分。（填序号。）}
\draw[line width=.7pt] (139.333,396.000)--(152.667,363.333)--(191.333,363.333)--(177.333,396.000)--cycle;
\draw[line width=.7pt] (245.333,395.333)--(260.667,364.667)--(275.333,395.333)--cycle;
\draw[line width=.7pt] (326.000,363.333)--(326.000,396.000)--(384.667,396.000)--cycle;
\draw[line width=.7pt] (416.667,394.000)--(432.667,364.667)--(466.667,364.667)--(482.667,394.000)--cycle;
\draw[line width=.7pt] (146.000,468.000)--(164.667,434.000)--(184.667,468.000)--cycle;
\draw[line width=.7pt] (238.667,434.667)--(282.667,434.667)--(282.667,468.000)--(238.667,468.000)--cycle;
\draw[line width=.7pt] (338.000,448.000)--(354.667,434.667)--(372.000,448.000)--(366.667,468.000)--(344.000,468.000)--cycle;
\draw[line width=.7pt] (439.333,467.333)--(444.000,434.000)--(460.667,434.000)--(456.000,467.333)--cycle;
\node[inner sep=0pt,align=center,font=\normalsize,] at (165.333,414.000) {①};
\node[inner sep=0pt,align=center,font=\normalsize,] at (260.667,408.667) {②};
\node[inner sep=0pt,align=center,font=\normalsize,] at (355.333,413.333) {③};
\node[inner sep=0pt,align=center,font=\normalsize,] at (450.000,408.667) {④};
\node[inner sep=0pt,align=center,font=\normalsize,] at (165.333,485.333) {⑤};
\node[inner sep=0pt,align=center,font=\normalsize,] at (260.667,485.333) {⑥};
\node[inner sep=0pt,align=center,font=\normalsize,] at (355.333,485.333) {⑦};
\node[inner sep=0pt,align=center,font=\normalsize,] at (450.000,481.333) {⑧};
\q{104}{769}{241}{12}{三角形：（\hspace{46pt}）}
\q{329}{769}{292}{12}{平行四边形：（\hspace{50pt}）}
\q{626}{769}{200}{12}{梯形：（\hspace{47pt}）}
\hdsection{4}{827}{在下面的点子图中画一画。}



\fill (75.333,673.333) circle (.7bp);
\fill (87.333,673.333) circle (.7bp);
\fill (99.333,673.333) circle (.7bp);
\fill (111.333,673.333) circle (.7bp);
\fill (123.333,673.333) circle (.7bp);
\fill (135.333,673.333) circle (.7bp);
\fill (147.333,673.333) circle (.7bp);
\fill (159.333,673.333) circle (.7bp);
\fill (171.333,673.333) circle (.7bp);
\fill (183.333,673.333) circle (.7bp);
\fill (195.333,673.333) circle (.7bp);
\fill (207.333,673.333) circle (.7bp);
\fill (219.333,673.333) circle (.7bp);
\fill (231.333,673.333) circle (.7bp);
\fill (243.333,673.333) circle (.7bp);
\fill (255.333,673.333) circle (.7bp);
\fill (267.333,673.333) circle (.7bp);
\fill (279.333,673.333) circle (.7bp);
\fill (291.333,673.333) circle (.7bp);
\fill (303.333,673.333) circle (.7bp);
\fill (315.333,673.333) circle (.7bp);
\fill (327.333,673.333) circle (.7bp);
\fill (339.333,673.333) circle (.7bp);
\fill (351.333,673.333) circle (.7bp);
\fill (363.333,673.333) circle (.7bp);
\fill (375.333,673.333) circle (.7bp);
\fill (387.333,673.333) circle (.7bp);
\fill (399.333,673.333) circle (.7bp);
\fill (411.333,673.333) circle (.7bp);
\fill (423.333,673.333) circle (.7bp);
\fill (435.333,673.333) circle (.7bp);
\fill (447.333,673.333) circle (.7bp);
\fill (459.333,673.333) circle (.7bp);
\fill (471.333,673.333) circle (.7bp);
\fill (483.333,673.333) circle (.7bp);
\fill (495.333,673.333) circle (.7bp);
\fill (507.333,673.333) circle (.7bp);
\fill (519.333,673.333) circle (.7bp);
\fill (531.333,673.333) circle (.7bp);
\fill (543.333,673.333) circle (.7bp);
\fill (75.333,685.333) circle (.7bp);
\fill (87.333,685.333) circle (.7bp);
\fill (99.333,685.333) circle (.7bp);
\fill (111.333,685.333) circle (.7bp);
\fill (123.333,685.333) circle (.7bp);
\fill (135.333,685.333) circle (.7bp);
\fill (147.333,685.333) circle (.7bp);
\fill (159.333,685.333) circle (.7bp);
\fill (171.333,685.333) circle (.7bp);
\fill (183.333,685.333) circle (.7bp);
\fill (195.333,685.333) circle (.7bp);
\fill (207.333,685.333) circle (.7bp);
\fill (219.333,685.333) circle (.7bp);
\fill (231.333,685.333) circle (.7bp);
\fill (243.333,685.333) circle (.7bp);
\fill (255.333,685.333) circle (.7bp);
\fill (267.333,685.333) circle (.7bp);
\fill (279.333,685.333) circle (.7bp);
\fill (291.333,685.333) circle (.7bp);
\fill (303.333,685.333) circle (.7bp);
\fill (315.333,685.333) circle (.7bp);
\fill (327.333,685.333) circle (.7bp);
\fill (339.333,685.333) circle (.7bp);
\fill (351.333,685.333) circle (.7bp);
\fill (363.333,685.333) circle (.7bp);
\fill (375.333,685.333) circle (.7bp);
\fill (387.333,685.333) circle (.7bp);
\fill (399.333,685.333) circle (.7bp);
\fill (411.333,685.333) circle (.7bp);
\fill (423.333,685.333) circle (.7bp);
\fill (435.333,685.333) circle (.7bp);
\fill (447.333,685.333) circle (.7bp);
\fill (459.333,685.333) circle (.7bp);
\fill (471.333,685.333) circle (.7bp);
\fill (483.333,685.333) circle (.7bp);
\fill (495.333,685.333) circle (.7bp);
\fill (507.333,685.333) circle (.7bp);
\fill (519.333,685.333) circle (.7bp);
\fill (531.333,685.333) circle (.7bp);
\fill (543.333,685.333) circle (.7bp);
\fill (75.333,697.333) circle (.7bp);
\fill (87.333,697.333) circle (.7bp);
\fill (99.333,697.333) circle (.7bp);
\fill (111.333,697.333) circle (.7bp);
\fill (123.333,697.333) circle (.7bp);
\fill (135.333,697.333) circle (.7bp);
\fill (147.333,697.333) circle (.7bp);
\fill (159.333,697.333) circle (.7bp);
\fill (171.333,697.333) circle (.7bp);
\fill (183.333,697.333) circle (.7bp);
\fill (195.333,697.333) circle (.7bp);
\fill (207.333,697.333) circle (.7bp);
\fill (219.333,697.333) circle (.7bp);
\fill (231.333,697.333) circle (.7bp);
\fill (243.333,697.333) circle (.7bp);
\fill (255.333,697.333) circle (.7bp);
\fill (267.333,697.333) circle (.7bp);
\fill (279.333,697.333) circle (.7bp);
\fill (291.333,697.333) circle (.7bp);
\fill (303.333,697.333) circle (.7bp);
\fill (315.333,697.333) circle (.7bp);
\fill (327.333,697.333) circle (.7bp);
\fill (339.333,697.333) circle (.7bp);
\fill (351.333,697.333) circle (.7bp);
\fill (363.333,697.333) circle (.7bp);
\fill (375.333,697.333) circle (.7bp);
\fill (387.333,697.333) circle (.7bp);
\fill (399.333,697.333) circle (.7bp);
\fill (411.333,697.333) circle (.7bp);
\fill (423.333,697.333) circle (.7bp);
\fill (435.333,697.333) circle (.7bp);
\fill (447.333,697.333) circle (.7bp);
\fill (459.333,697.333) circle (.7bp);
\fill (471.333,697.333) circle (.7bp);
\fill (483.333,697.333) circle (.7bp);
\fill (495.333,697.333) circle (.7bp);
\fill (507.333,697.333) circle (.7bp);
\fill (519.333,697.333) circle (.7bp);
\fill (531.333,697.333) circle (.7bp);
\fill (543.333,697.333) circle (.7bp);
\fill (75.333,709.333) circle (.7bp);
\fill (87.333,709.333) circle (.7bp);
\fill (99.333,709.333) circle (.7bp);
\fill (111.333,709.333) circle (.7bp);
\fill (123.333,709.333) circle (.7bp);
\fill (135.333,709.333) circle (.7bp);
\fill (147.333,709.333) circle (.7bp);
\fill (159.333,709.333) circle (.7bp);
\fill (171.333,709.333) circle (.7bp);
\fill (183.333,709.333) circle (.7bp);
\fill (195.333,709.333) circle (.7bp);
\fill (207.333,709.333) circle (.7bp);
\fill (219.333,709.333) circle (.7bp);
\fill (231.333,709.333) circle (.7bp);
\fill (243.333,709.333) circle (.7bp);
\fill (255.333,709.333) circle (.7bp);
\fill (267.333,709.333) circle (.7bp);
\fill (279.333,709.333) circle (.7bp);
\fill (291.333,709.333) circle (.7bp);
\fill (303.333,709.333) circle (.7bp);
\fill (315.333,709.333) circle (.7bp);
\fill (327.333,709.333) circle (.7bp);
\fill (339.333,709.333) circle (.7bp);
\fill (351.333,709.333) circle (.7bp);
\fill (363.333,709.333) circle (.7bp);
\fill (375.333,709.333) circle (.7bp);
\fill (387.333,709.333) circle (.7bp);
\fill (399.333,709.333) circle (.7bp);
\fill (411.333,709.333) circle (.7bp);
\fill (423.333,709.333) circle (.7bp);
\fill (435.333,709.333) circle (.7bp);
\fill (447.333,709.333) circle (.7bp);
\fill (459.333,709.333) circle (.7bp);
\fill (471.333,709.333) circle (.7bp);
\fill (483.333,709.333) circle (.7bp);
\fill (495.333,709.333) circle (.7bp);
\fill (507.333,709.333) circle (.7bp);
\fill (519.333,709.333) circle (.7bp);
\fill (531.333,709.333) circle (.7bp);
\fill (543.333,709.333) circle (.7bp);
\fill (75.333,721.333) circle (.7bp);
\fill (87.333,721.333) circle (.7bp);
\fill (99.333,721.333) circle (.7bp);
\fill (111.333,721.333) circle (.7bp);
\fill (123.333,721.333) circle (.7bp);
\fill (135.333,721.333) circle (.7bp);
\fill (147.333,721.333) circle (.7bp);
\fill (159.333,721.333) circle (.7bp);
\fill (171.333,721.333) circle (.7bp);
\fill (183.333,721.333) circle (.7bp);
\fill (195.333,721.333) circle (.7bp);
\fill (207.333,721.333) circle (.7bp);
\fill (219.333,721.333) circle (.7bp);
\fill (231.333,721.333) circle (.7bp);
\fill (243.333,721.333) circle (.7bp);
\fill (255.333,721.333) circle (.7bp);
\fill (267.333,721.333) circle (.7bp);
\fill (279.333,721.333) circle (.7bp);
\fill (291.333,721.333) circle (.7bp);
\fill (303.333,721.333) circle (.7bp);
\fill (315.333,721.333) circle (.7bp);
\fill (327.333,721.333) circle (.7bp);
\fill (339.333,721.333) circle (.7bp);
\fill (351.333,721.333) circle (.7bp);
\fill (363.333,721.333) circle (.7bp);
\fill (375.333,721.333) circle (.7bp);
\fill (387.333,721.333) circle (.7bp);
\fill (399.333,721.333) circle (.7bp);
\fill (411.333,721.333) circle (.7bp);
\fill (423.333,721.333) circle (.7bp);
\fill (435.333,721.333) circle (.7bp);
\fill (447.333,721.333) circle (.7bp);
\fill (459.333,721.333) circle (.7bp);
\fill (471.333,721.333) circle (.7bp);
\fill (483.333,721.333) circle (.7bp);
\fill (495.333,721.333) circle (.7bp);
\fill (507.333,721.333) circle (.7bp);
\fill (519.333,721.333) circle (.7bp);
\fill (531.333,721.333) circle (.7bp);
\fill (543.333,721.333) circle (.7bp);
\fill (75.333,733.333) circle (.7bp);
\fill (87.333,733.333) circle (.7bp);
\fill (99.333,733.333) circle (.7bp);
\fill (111.333,733.333) circle (.7bp);
\fill (123.333,733.333) circle (.7bp);
\fill (135.333,733.333) circle (.7bp);
\fill (147.333,733.333) circle (.7bp);
\fill (159.333,733.333) circle (.7bp);
\fill (171.333,733.333) circle (.7bp);
\fill (183.333,733.333) circle (.7bp);
\fill (195.333,733.333) circle (.7bp);
\fill (207.333,733.333) circle (.7bp);
\fill (219.333,733.333) circle (.7bp);
\fill (231.333,733.333) circle (.7bp);
\fill (243.333,733.333) circle (.7bp);
\fill (255.333,733.333) circle (.7bp);
\fill (267.333,733.333) circle (.7bp);
\fill (279.333,733.333) circle (.7bp);
\fill (291.333,733.333) circle (.7bp);
\fill (303.333,733.333) circle (.7bp);
\fill (315.333,733.333) circle (.7bp);
\fill (327.333,733.333) circle (.7bp);
\fill (339.333,733.333) circle (.7bp);
\fill (351.333,733.333) circle (.7bp);
\fill (363.333,733.333) circle (.7bp);
\fill (375.333,733.333) circle (.7bp);
\fill (387.333,733.333) circle (.7bp);
\fill (399.333,733.333) circle (.7bp);
\fill (411.333,733.333) circle (.7bp);
\fill (423.333,733.333) circle (.7bp);
\fill (435.333,733.333) circle (.7bp);
\fill (447.333,733.333) circle (.7bp);
\fill (459.333,733.333) circle (.7bp);
\fill (471.333,733.333) circle (.7bp);
\fill (483.333,733.333) circle (.7bp);
\fill (495.333,733.333) circle (.7bp);
\fill (507.333,733.333) circle (.7bp);
\fill (519.333,733.333) circle (.7bp);
\fill (531.333,733.333) circle (.7bp);
\fill (543.333,733.333) circle (.7bp);
\fill (75.333,745.333) circle (.7bp);
\fill (87.333,745.333) circle (.7bp);
\fill (99.333,745.333) circle (.7bp);
\fill (111.333,745.333) circle (.7bp);
\fill (123.333,745.333) circle (.7bp);
\fill (135.333,745.333) circle (.7bp);
\fill (147.333,745.333) circle (.7bp);
\fill (159.333,745.333) circle (.7bp);
\fill (171.333,745.333) circle (.7bp);
\fill (183.333,745.333) circle (.7bp);
\fill (195.333,745.333) circle (.7bp);
\fill (207.333,745.333) circle (.7bp);
\fill (219.333,745.333) circle (.7bp);
\fill (231.333,745.333) circle (.7bp);
\fill (243.333,745.333) circle (.7bp);
\fill (255.333,745.333) circle (.7bp);
\fill (267.333,745.333) circle (.7bp);
\fill (279.333,745.333) circle (.7bp);
\fill (291.333,745.333) circle (.7bp);
\fill (303.333,745.333) circle (.7bp);
\fill (315.333,745.333) circle (.7bp);
\fill (327.333,745.333) circle (.7bp);
\fill (339.333,745.333) circle (.7bp);
\fill (351.333,745.333) circle (.7bp);
\fill (363.333,745.333) circle (.7bp);
\fill (375.333,745.333) circle (.7bp);
\fill (387.333,745.333) circle (.7bp);
\fill (399.333,745.333) circle (.7bp);
\fill (411.333,745.333) circle (.7bp);
\fill (423.333,745.333) circle (.7bp);
\fill (435.333,745.333) circle (.7bp);
\fill (447.333,745.333) circle (.7bp);
\fill (459.333,745.333) circle (.7bp);
\fill (471.333,745.333) circle (.7bp);
\fill (483.333,745.333) circle (.7bp);
\fill (495.333,745.333) circle (.7bp);
\fill (507.333,745.333) circle (.7bp);
\fill (519.333,745.333) circle (.7bp);
\fill (531.333,745.333) circle (.7bp);
\fill (543.333,745.333) circle (.7bp);
\fill (75.333,757.333) circle (.7bp);
\fill (87.333,757.333) circle (.7bp);
\fill (99.333,757.333) circle (.7bp);
\fill (111.333,757.333) circle (.7bp);
\fill (123.333,757.333) circle (.7bp);
\fill (135.333,757.333) circle (.7bp);
\fill (147.333,757.333) circle (.7bp);
\fill (159.333,757.333) circle (.7bp);
\fill (171.333,757.333) circle (.7bp);
\fill (183.333,757.333) circle (.7bp);
\fill (195.333,757.333) circle (.7bp);
\fill (207.333,757.333) circle (.7bp);
\fill (219.333,757.333) circle (.7bp);
\fill (231.333,757.333) circle (.7bp);
\fill (243.333,757.333) circle (.7bp);
\fill (255.333,757.333) circle (.7bp);
\fill (267.333,757.333) circle (.7bp);
\fill (279.333,757.333) circle (.7bp);
\fill (291.333,757.333) circle (.7bp);
\fill (303.333,757.333) circle (.7bp);
\fill (315.333,757.333) circle (.7bp);
\fill (327.333,757.333) circle (.7bp);
\fill (339.333,757.333) circle (.7bp);
\fill (351.333,757.333) circle (.7bp);
\fill (363.333,757.333) circle (.7bp);
\fill (375.333,757.333) circle (.7bp);
\fill (387.333,757.333) circle (.7bp);
\fill (399.333,757.333) circle (.7bp);
\fill (411.333,757.333) circle (.7bp);
\fill (423.333,757.333) circle (.7bp);
\fill (435.333,757.333) circle (.7bp);
\fill (447.333,757.333) circle (.7bp);
\fill (459.333,757.333) circle (.7bp);
\fill (471.333,757.333) circle (.7bp);
\fill (483.333,757.333) circle (.7bp);
\fill (495.333,757.333) circle (.7bp);
\fill (507.333,757.333) circle (.7bp);
\fill (519.333,757.333) circle (.7bp);
\fill (531.333,757.333) circle (.7bp);
\fill (543.333,757.333) circle (.7bp);
\fill (75.333,769.333) circle (.7bp);
\fill (87.333,769.333) circle (.7bp);
\fill (99.333,769.333) circle (.7bp);
\fill (111.333,769.333) circle (.7bp);
\fill (123.333,769.333) circle (.7bp);
\fill (135.333,769.333) circle (.7bp);
\fill (147.333,769.333) circle (.7bp);
\fill (159.333,769.333) circle (.7bp);
\fill (171.333,769.333) circle (.7bp);
\fill (183.333,769.333) circle (.7bp);
\fill (195.333,769.333) circle (.7bp);
\fill (207.333,769.333) circle (.7bp);
\fill (219.333,769.333) circle (.7bp);
\fill (231.333,769.333) circle (.7bp);
\fill (243.333,769.333) circle (.7bp);
\fill (255.333,769.333) circle (.7bp);
\fill (267.333,769.333) circle (.7bp);
\fill (279.333,769.333) circle (.7bp);
\fill (291.333,769.333) circle (.7bp);
\fill (303.333,769.333) circle (.7bp);
\fill (315.333,769.333) circle (.7bp);
\fill (327.333,769.333) circle (.7bp);
\fill (339.333,769.333) circle (.7bp);
\fill (351.333,769.333) circle (.7bp);
\fill (363.333,769.333) circle (.7bp);
\fill (375.333,769.333) circle (.7bp);
\fill (387.333,769.333) circle (.7bp);
\fill (399.333,769.333) circle (.7bp);
\fill (411.333,769.333) circle (.7bp);
\fill (423.333,769.333) circle (.7bp);
\fill (435.333,769.333) circle (.7bp);
\fill (447.333,769.333) circle (.7bp);
\fill (459.333,769.333) circle (.7bp);
\fill (471.333,769.333) circle (.7bp);
\fill (483.333,769.333) circle (.7bp);
\fill (495.333,769.333) circle (.7bp);
\fill (507.333,769.333) circle (.7bp);
\fill (519.333,769.333) circle (.7bp);
\fill (531.333,769.333) circle (.7bp);
\fill (543.333,769.333) circle (.7bp);

\begin{scope}[color=answerred,draw=answerred,text=answerred]% 高清文字版专用填空位置；原页答案另存于 ../overlays。
\ans{268.667}{186.667}{40.000}{14}{$\checkmark$}
\ans{352.000}{186.667}{40.000}{14}{$\checkmark$}
\ans{162.000}{263.333}{32.000}{12}{$5$}
\ans{162.000}{292.667}{32.000}{12}{$3$}
\ans{420.667}{263.333}{32.000}{12}{$3$}
\ans{420.667}{292.667}{32.000}{12}{$1$}
\ans{129.613}{512.343}{66.000}{11}{②③⑤}
\ans{305.517}{512.343}{70.000}{11}{①⑥⑧}
\ans{477.115}{512.343}{43.333}{11}{④}
\begin{scope}[scale=0.6666667]\draw[answerred,line width=1pt] (149,1136)--(149,1064)--(221,1136)--cycle;\draw (149,1125)--(160,1125)--(160,1136);\draw (144,1100)--(154,1100) (185,1131)--(185,1141);\node[answerred,above] at (185,1050) {①};\draw[answerred,line width=1pt] (365,1136)--(401,1064)--(473,1064)--(509,1136)--cycle;\draw[answerred,dashed] (401,1064)--(401,1136);\draw (401,1125)--(412,1125)--(412,1136);\node[above] at (437,1061) {上底};\node[below] at (437,1136) {下底};\node[left] at (401,1100) {高};\node[above] at (365,1050) {②};\draw[answerred,line width=1pt] (617,1136)--(653,1064)--(761,1064)--(725,1136)--cycle;\draw[answerred,dashed] (653,1064)--(653,1136);\draw (653,1125)--(664,1125)--(664,1136);\node[left] at (653,1100) {高};\node[above] at (616,1050) {③};\end{scope}
\end{scope}
\end{scope}\pgfresetboundingbox\path[use as bounding box] (69.0,667.0) rectangle (549.0,776.0);
\end{tikzpicture}}\par

\PlainItem{① 画出有一个角是直角的等腰三角形。}
\par

\PlainItem{② 画一个梯形，并画出它的一条高，标出上底和下底。}
\PlainItem{③ 画一个平行四边形，并画出它的一条高。}




\QHeading{第4题　开放作图的其他答案}

\Answer{等腰直角三角形可改变两条直角边的共同长度；梯形可改变上下底和高；平行四边形可改变底、高和倾斜方向。三角形指定底的高仍只从对面顶点作出；平行四边形和梯形可从另一条平行边的不同点作垂线段。新增高用虚线，必要的底边延长线也用虚线。}

\Answer{作高复核：高用虚线，旁边写“高”，垂足处画直角符号；以题目指定的底为准，必要时用虚线延长底边。}
